\chapter*{Razširjeni povzetek}
\addcontentsline{toc}{chapter}{Razširjeni povzetek}

\TODO{
To poglavje napiši na koncu, ko bodo znani rezultati. Struktura naj bo podobna raziskovalnemu članku:
\begin{itemize}
    \item novo izhodišče: glavna evalvacija naj bo binarna klasifikacija in primerjava modelov pri množicah \emph{samo pogled}, \emph{pogled in zenici} in \emph{vsi signali};
    \item \GazeMAE{} opiši kot referenco za množico \emph{samo pogled}; \emph{samo zenici}/MOMENT omeni samo, če bo eksperiment dokončan;
    \item motivacija: zakaj grafovska predstavitev podatkov sledilnika pogleda;
    \item predlagana metoda: kako iz časovnih oken nastanejo grafi;
    \item eksperimentalni protokol: \MAHNOB{}, 3-razredna valenca/vzburjenost, LOO po subjektih in LOO po posnetkih;
    \item rezultati: primerjava s primerjalnimi modeli, \GazeMAE{} in GNN arhitekturami;
    \item sklep: kaj grafovska predstavitev prispeva in kje so omejitve.
\end{itemize}
}

\section*{I Grafovske nevronske mreže in motivacija}

\TODO{
Kratko povzemi, kaj so grafovske nevronske mreže, zakaj so primerne za podatke z relacijsko, prostorsko ali časovno strukturo in zakaj je okno meritev sledilnika pogleda smiselno obravnavati kot prostorsko-časovni graf.
}

\section*{II Grafovska predstavitev podatkov sledilnika pogleda}

\TODO{
\begin{itemize}
    \item Razširjeni povzetek naj grafovsko predstavitev samo povzame; natančne definicije ostanejo v poglavju~\ref{sec:grafovska-predstavitev}.
    \item Ne ponavljaj formul za konstrukcijo povezav, značilk povezav ali normalizacijo uteži.
    \item Uteži povezav omeni samo kot arhitekturno komponento, definirano v poglavju~\ref{sec:modeli}; ne kot del osnovne definicije vhodnega grafa.
    \item vozlišče = ena meritev sledilnika pogleda;
    \item značilke vozlišča = koordinati pogleda $(x,~y)$, leva/desna zenica, \texttt{time-window-normalized}, \texttt{distance-avg}, \texttt{fixation-duration};
    \item časovne povezave = bližina vzorcev v času, $k_t=2$;
    \item prostorske povezave = bližina vzorcev v prostoru pogleda, $k_s=2$;
    \item fiksacijske povezave = razširjene povezave znotraj fiksacije, $k_f=3$;
    \item uteži povezav = naučene predznačene uteži iz značilk relacije;
    \item \texttt{fixation-index} = grupiranje, ne numerična značilka vozlišča.
\end{itemize}
}

\section*{III Modeliranje in eksperimentalna evalvacija}

\TODO{
Kratko povzemi uporabljeno podatkovno zbirko \MAHNOB{}, nalogi 3-razredne valence in 3-razredne vzburjenosti, validacijo LOO po subjektih in LOO po posnetkih ter primerjane modele: SVM, LightGBM, MLP, \GazeMAE{} + MLP, osnovni GCN in predlagani GNN.
}

\section*{IV Rezultati in sklep}

\TODO{
Dopolni po zaključku eksperimentov: kateri model je najboljši pri valenci, kateri pri vzburjenosti, ali GNN izboljša rezultate glede na ne-grafovske primerjalne modele, katere komponente grafa najbolj prispevajo k rezultatu ter glavne omejitve in prihodnje delo.
}
